
% how to change Contents to Table of Contents
\addto\captionsenglish{% Replace "english" with the language you use
  \renewcommand{\contentsname}%
    {Table of Contents}%
}

% to change the name of Abbreviations to Acronyms
% not needed if use use entry types and define those
% \renewcommand{\abbreviationsname}{Acronyms}

% to allow line comments in algorithms
\algnewcommand{\LineComment}[1]{\State \(\triangleright\) #1}

% to declare abs and norm
\DeclarePairedDelimiter\abs{\lvert}{\rvert}%
\DeclarePairedDelimiter\norm{\lVert}{\rVert}%

% Swap the definition of \abs* and \norm*, so that \abs
% and \norm resizes the size of the brackets, and the 
% starred version does not.
\makeatletter
\let\oldabs\abs
\def\abs{\@ifstar{\oldabs}{\oldabs*}}
%
\let\oldnorm\norm
\def\norm{\@ifstar{\oldnorm}{\oldnorm*}}
\makeatother

% heatmap cell colors, extracted from the original figure (images/ch4_generalization.png).
% Rule: each panel normalizes MAE linearly to t in [0,1] (NBD: min 0.023, max 0.668;
% NRD: min 1.7, max 28.2), then maps t onto a green -> yellow -> red tint scale
% (pure hue blended with white). Greens/reds match xcolor blends exactly;
% yellows use the figure's exact (slightly warm) hex values.
% Full per-value mapping and update procedure: 修改/图4.1配色规律.md
\definecolor{hg10}{HTML}{E5FFE5} % green!10!white
\definecolor{hg15}{HTML}{D8FFD8} % green!15!white
\definecolor{hg20}{HTML}{CCFFCC} % green!20!white
\definecolor{hg25}{HTML}{BFFFBF} % green!25!white
\definecolor{hg30}{HTML}{B2FFB2} % green!30!white
\definecolor{hy08}{HTML}{FEFDE9} % yellow!8!white (figure: slightly warm)
\definecolor{hy12}{HTML}{FEFCDF} % yellow!12!white
\definecolor{hy16}{HTML}{FEFCD6} % yellow!16!white
\definecolor{hy20}{HTML}{FEFBCC} % yellow!20!white
\definecolor{hy24}{HTML}{FEFAC2} % yellow!24!white
\definecolor{hy32}{HTML}{FEF8AD} % yellow!32!white
\definecolor{hy36}{HTML}{FEF8A3} % yellow!36!white
\definecolor{hr50}{HTML}{FF7F7F} % red!50!white
\definecolor{hr55}{HTML}{FF7272} % red!55!white
\definecolor{hr60}{HTML}{FF6666} % red!60!white
\definecolor{hr70}{HTML}{FF4C4C} % red!70!white
\definecolor{hr75}{HTML}{FF3F3F} % red!75!white
\definecolor{hr80}{HTML}{FF3333} % red!80!white
\newcommand{\hc}[2]{\cellcolor{#1}#2} % colored heatmap cell: #1=color name, #2=text
